%!TEX encoding = UTF8
%!TEX root = 0-notes.tex

\unchapter{Calcul littéral}
\fancyhead[C]{\textbf{Chapitre \thechapter~— Calcul littéral}}
\addcontentsline{toc}{section}{Méthodologie}

\axiome{
	Pour trois nombres réels $x, y, z \in\R$, nous avons
	\begin{enumerate}
		\item
		La distributivité de $\times$ sur $+$ :
			\[ x\times(y+z) = x\times y + x\times z. \]
		\item
		L'associativité de l'addition :
			\[ x+(y+z) = (x+y)+z = x+y+z.\]
		\item
		L'associativité de la multiplication :
			\[ x\times(y\times z) = (x\times y)\times z = x\times y\times z.\]
	\end{enumerate}
}{}

\nt{
	Les axiomes s'appliquent à la division et la soustraction car
	\begin{center}
	« Diviser c'est multiplier par l'inverse. »
	
	« Soustraire c'est ajouter l'opposé. »
	\end{center}
}

\nt{
	Lorsqu'une opération est omise, c'est une multiplication.
	La multiplication est prioritaire sur l'addition. Les parenthèses isolent des expressions.
}

\ex{}{
		\[ 
		\underbrace{3 (x + 4) + 2}_{\emphindex{forme factorisée}}  
		= 
		\underbrace{3x + 3\times4 + 2}_ {\emphindex{forme développée}}  
		= 
		\underbrace{3x + 14}_{\emphindex{forme développée réduite}}
		. \]
		\[ 
		\underbrace{xy + 2yz}_{\emphindex{forme développée réduite}}
		= 
		\underbrace{y(x+2z)}_{\emphindex{forme factorisée}}
		. \]
}{}

\dfn{puissance sur $\N^*$}{
	Soit $n \in \N^*$ un entier naturel non nul.
	Alors \og $x$ puissance $n$ \fg~est égal à
		\[ x^{n} = \underbrace{x \times x \times \cdots \times x}_{\text{$n$ fois}}. \]
	En particulier,
		\begin{align*}
			x^1 = x && x^2 = x\times x && x^3 = x \times x \times x
		\end{align*}
}{dfn:puissance-N}

\nt{
	La puissance 0 est égale à 1 par convention  : $x^0 = 1$.
}

\thm{}{
	Soient $a, b, c, d\in\R$ quatre nombres réels.
	Alors
	\begin{align*}
		&(a+b)(c+d) = (a)(c) + (a)(d) + (b)(c) + (b)(d) &\text{(double distributivité)} \\
	&\left.
	\begin{array}{cl}
		(a+b)(a-b) &= (a)^2 - (b)^2 \\
		(a+b)^2 &= (a)^2 + (b)^2 + 2(a)(b)\\
		(a-b)^2 &= (a)^2 + (b)^2 - 2(a)(b)
	\end{array}\right\}  &\text{(identités remarquables)} 
	\end{align*}
	%$$
}{}

\prop{propriétés des puissances}{
	Soient $x, y\in\R$ et  $m, n \in \N$. Alors
		\begin{align*}
			x^{m} \times x^{n} = x^{m+n},  && \bigl( x^{m}\bigr)^n = x^{m\times n}, && \et && (xy)^n = x^n y^n.
		\end{align*}
%		\[ x^{m} \times x^{n} = x^{m+n}, \]
%	et 
%		\[ \bigl( x^{m}\bigr)^n = x^{m\times n}. \]
%	En particulier, on a 
%		\begin{align*}
%			q^0 = 1 && q^{-1} = \dfrac1q &&	q^{-a} = \dfrac1{q^a}
%		\end{align*}
}{prop:propriétés-puissances}

\ex{}{
		\[ 
		3(2x-1)^2 + x^2
		= 
		3\underbrace{\left[(2x)^2+(1)^2 - 2(2x)(1)\right]}_{(2x-1)^2} + x^2
		= 
		3\bigl[4x^2+1^2 - 4x\bigr] + x^2
		= 
		\bigl[12x^2+3 - 12x\bigr] + x^2
		= 
		\underbrace{13x^2}_{12x^2 + x^2}+3 - 12x
		. \]
}{}

\newpage
\addcontentsline{toc}{section}{Exercices}
\fancyhead[C]{\textbf{Chapitre \thechapter~— Exercices}}


\exe{}{
	Développer et réduire les expressions suivantes pour obtenir une expression de la forme
		\[ ax^2 + bx + c, \]
	où $a, b, c \in\R$ sont des nombres réels (possiblement nuls).
	
	\begin{multicols}{2}
	\begin{enumerate}
		\item $x + 2x + 3x$
		\item $(1+3x) - (8 - 2x)$
		\item $x - 2(2-x)$
		\item $x(3-x) + 8(x+1)$
		\item $(x+3)(x-1) - 2x + 1$
		\item $(x+1)^2 + 4x - 4$
		\item $(3x-1)^2 - 9x^2 - 4$
		\item $x(3x-1)  + 3x^2 +12$
	\end{enumerate}
	\end{multicols}
}{exe:dev-red}{}

\exe{}{
	Développer les expressions algébriques suivantes pour obtenir une expression de la forme
		\[ ax^2 + bx + c, \]
	où $a, b, c \in\R$ sont des nombres réels (possiblement nuls).
		\begin{multicols}{2}
		\begin{enumerate}
			\item $f(x) = (1+x)^2$
			\item $g(x) = (x-3)^2$
			\item $h(x) = (3-x)^2$
			\item $F(x) = (3 + 2x)^2$
			\item $G(x) = (3x - 7)^2$
			\item $H(x) = (-7x - 2)^2$
		\end{enumerate}
		\end{multicols}
}{exe:id-rem}{
	\begin{align*}
		f(x) &= (1+x)^2 = 1^2 + 2(1)(x) + x^2 = 1 + 2x + x^2 \\
		g(x) &= (x-3)^2 = x^2 - 2(x)(3) + 3^2 = x^2 - 6x + 9 \\
		h(x) &= (3-x)^2 = 3^2 - 2(3)(x) + x^2 = 9 - 6x + x^2 \\
		F(x) &= (3 + 2x)^2 = 3^2 + 2(3)(2x) + (2x)^2 = 9 + 12x + 4x^2 \\
		G(x) &= (3x - 7)^2 = (3x)^2 - 2(3x)(7) + 7^2 = 9x^2 - 42x + 49 \\
		H(x) &= (-7x - 2)^2 = (7x + 2)^2 = (7x)^2 + 2(7x)(2) + 2^2 = 49x^2 + 28x + 4
	\end{align*}

}



\exe{}{
	Développer les expressions suivantes.
		\begin{multicols}{3}
		\begin{enumerate}%[label=\arabic*)]
			\item $\left(1+\sqrt2\right)\left(1-\sqrt2\right)$
			\item $\left(1+\sqrt2\right)^2$
			\item $\left(1-\sqrt2\right)^2$
			\item $\left(\sqrt{3} - 2\sqrt2\right)^2$
			\item $\left(\sqrt5x + 2 \right)^2$
			\item $\left(\sqrt5x - \sqrt7\right)^2$
		\end{enumerate}
		\end{multicols}
}{exe:développement-sqrt}{
	\begin{enumerate}
		\item $\left(1+\sqrt2\right)\left(1-\sqrt2\right) = 1-\sqrt2^2 = 1-2 = -1$
		\item $\left(1+\sqrt2\right)^2 = 1 + \sqrt2^2 + 2\sqrt2 = 3 + 2\sqrt2$
		\item $\left(1-\sqrt2\right)^2 = 1 + 2 - 2\sqrt2 = 3 - 2\sqrt2$
		\item $\left(\sqrt{3} - 2\sqrt2\right)^2 = 3 + 4\times2 - 2\times\sqrt3\times2\times\sqrt2 = 11 - 4\sqrt6$
		\item $\left(\sqrt5x + 2 \right)^2 = 5x^2 + 4 + 4\sqrt5x$
		\item $\left(\sqrt5x - \sqrt7\right)^2 = 5x^2 + 7 - 2\sqrt{35}x$
	\end{enumerate}
}

\exe{, difficulty=1}{
	En étant donné que  $\cos^2(x) + \sin^2(x) = 1$ pour tout $x\in\R$ et que
		\begin{align*}
			\sin(15\degree) = \dfrac{\sqrt{6}-\sqrt{2}}4 && \text{ et } && \cos(36\degree) = \dfrac{1+\sqrt{5}}4,
		\end{align*}
	vérifier que $\cos(15\degree) = \frac{\sqrt{6}+\sqrt{2}}4$ et donner la valeur \textbf{exacte} de $\sin(36\degree)$.
}{exe:c2-s2-1}{
	Premièrement, vérifions l'égalité $\sin^2(15^\circ) + \cos^2(15^\circ) = 1$ avec les valeurs fournies.
	\begin{align*}
		\left(\frac{\sqrt{6}+\sqrt{2}}4\right)^2 + \left(\dfrac{\sqrt{6}-\sqrt{2}}4\right)^2 
		&= \frac{6+2+2\sqrt{12}}{16} + \frac{6+2-2\sqrt{12}}{16} \\
		&= \frac{8+8}{16} = 1
	\end{align*}
	Or, pour un $y$ donnée, il y a deux solutions $x$ à l'équation $x^2 + y^2 = 1$ : une positive, et l'autre négative.
	Comme $\frac{\sqrt{6}+\sqrt{2}}4 > 0$, on a trouvé la solution positive, bien égale à $\cos(15^\circ)$ car celui-ci est positif d'après le cours.
	
	Deuxièmement, $\sin(36^\circ)$ est la solution $x\in\R$ positive de
		\begin{align*}
			x^2 + \cos^2(36^\circ) &= 1 \\
			x^2 &= 1 - \left(\dfrac{1+\sqrt{5}}4\right)^2 \\
			x^2 &= 1 - \dfrac{1+5+2\sqrt5}{16} \\
			x^2 &= \dfrac{10 - 2 \sqrt5}{16} = \dfrac{5-\sqrt5}8
		\end{align*}
	Par conséquent, $\sin(36^\circ) = \sqrt{\frac{5-\sqrt5}8}$.
}


\exe{, difficulty=2}{
	Supposons que le polynôme du second degré $p(x) = ax^2 + bx + c$ s'écrive de la forme $p(x) = a'(x-\alpha)^2 - \beta$.
	Montrer, par identification des coefficients en $x^2, x, 1$ que $a'=a$, puis que $\alpha = \frac{-b}{2a}$, et enfin que $\beta = \frac{b^2 - 4ac}{4a}$.
	
	Écrire $p(x) = 3x^2 - 3x - 1$ sous la forme $a(x-\alpha)^2 - \beta$ puis déduire le minimum de $p$ ainsi que l'antécédent le réalisant.
	
	\emph{La quantité $b^2 - 4ac$ est le \emphindex{discriminant} du polynôme et est notée $\Delta$ (lu « delta »).}
}{exe:déterminant-2nddeg}{
%	À rendre pour possibilité de points bonus à la prochaine évaluation.
%	Ni correction ni points ne seront attribués à un travail suspect ou ne démontrant pas une bonne compréhension de l'exercice.
%	
%	Le lemme suivant est (sans doute) nécessaire. Il peut être admis ou démontré.
%	
%	\begin{lemma}
%		Soient $a, b, c, a', b', c'$ six réels. Si l'identité
%			\[ ax^2 + bx + c = a'x^2 + b'x + c' \]
%		tient pour tout $x\in\R$, alors $a=a', b=b',$ et $c=c'$.
%	\end{lemma}
%	
%	Pour démontrer le lemme, vous pouvez d'abord démontrer le résultat suivant (qui est en fait équivalent).
%	
%	\begin{lemma}
%		Soient $a, b, c$ trois réels. Si l'identité
%			\[ ax^2 + bx + c = 0 \]
%		tient pour tout $x\in\R$, alors $a=b=c=0$.
%		
%		On dit alors que $1, x,$ et $x^2$ sont \emph{linéairement indépendants}.
%	\end{lemma}
%	
%	
	Forçons l'égalité $p(x) = a'(x-\alpha)^2 - \beta$ et voyons ce que chaque coefficient nous impose.
		\begin{align*}
			ax^2 + bx + c &= a'(x-\alpha)^2 - \beta \\
							&= a'(x^2 + \alpha^2 - 2\alpha x) - \beta \\
							&= a'x^2 + (-2a'\alpha)x + (a'\alpha^2 - \beta)
		\end{align*}
	 Les coefficients en $x^2$ doivent être égaux, donc $a' = a$ est forcé.
	 L'équation devient donc
	 	\[ ax^2 + bx + c = ax^2 + (-2a\alpha)x + (a\alpha^2-\beta). \]
	Les coefficients en $x$ doivent être égaux, donc $b = -2a\alpha$ est forcé.
	Il suit donc que
		\[ \alpha = \dfrac{-b}{2a}. \]
	L'équation devient donc
		\[ ax^2 + bx + c = ax^2 + bx + \left(\dfrac{b^2}{4a} - \beta\right). \]
	Le coefficients constants (en 1) doivent être égaux, donc $c = \dfrac{b^2}{4a} - \beta$.
	Il suit donc que
		\[ \beta = \dfrac{b^2}{4a} - c = \dfrac{b^2 - 4ac}{4a}, \]
	comme requis.
	
	Dans le cas de $p(x) = 3x^2 - 3x - 1$, on spécialise les formules pour $\alpha$ et $\beta$ lorsque $a=3, b=-3,$ et $c=-1$ pour trouver
		\begin{align*}
			\alpha = \dfrac{-(-3)}{2\times3} = \dfrac12 && \et && \beta = \dfrac{(-3)^2 - 4\times3\times(-1)}{4\times3} = \dfrac{21}{12} = \dfrac74.
		\end{align*}
	Par conséquent,
		\[ 3x^2  - 3x -1 = 3\left(x - \dfrac12\right)^2 - \dfrac74, \]
	égalité qu'on peut vérifier en redéveloppant le membre de droite pour se rassurer :
		\begin{align*}
			3\left(x - \dfrac12\right)^2 - \dfrac74 &= 3\left[ x^2 + \dfrac14 - x\right] - \dfrac74, \\
												&= 3x^2 + \dfrac34 - 3x - \dfrac74, \\
												&= 3x^2 - 3x - 1 = p(x).
		\end{align*}
	En conclusion, $p$ atteint son minimum $-\frac74$ en $x^\star = \frac12$.
}


\exe{, difficulty=2}{
	On appelle \emph{nombre de Mersenne}\footnotemark 
	un nombre prenant la forme $2^n - 1$ avec $n\in\N$.
	\begin{enumerate}
		\item
		Donner les 15 plus petits nombres de Mersenne.
		\item
		Parmis ces nombres, lesquels sont premiers ? L'utilisation d'un ordinateur est pertinente.
		\item
		Que dire des exposants $n$ lorsque $2^n -1$ est premier ?
		\item
		Montrer que pour tout $q\in\R$ et $d\geq1$ entier,
			\[ (q-1)\left(1+q+q^2 + q^3 + \cdots + q^{d-1}\right) = q^d - 1. \]
		\item
		En déduire que si $n$ n'est pas premier, alors $2^n -1$ n'est pas premier non plus. 
	\end{enumerate}
}{exe:Mersenne}{
	%À rendre pour possibilité de points bonus à la prochaine évaluation.
	%Ni correction ni points ne seront attribués à un travail suspect ou ne démontrant pas une bonne compréhension de l'exercice.
	
	%Voici un programme Python permettant de trouver les diviseurs d'un nombre.
	%\python{diviseurs}
}
\footnotetext{Marin Mersenne (1588-1648), physicien, mathématicien, musicologue et philosophe français.}

\exe{, difficulty=2}{
	Ce exercice traite du \emph{théorème des deux carrés} de Fermat\footnotemark.
	On dit que $n\in\N$ est somme de deux carrés si $n=a^2 + b^2$ avec $a, b \in \N$ entiers.
	\begin{enumerate}
		\item 
		Montrer que 2 et 4 sont somme de deux carrés, mais que 3 et 7 ne le sont pas.
		%\item 
		%Montrer qu'un nombre de la forme $4k+3$ où $k\in\N$ ne peut pas être somme de deux carrés.
		%
		%\emph{Indication : justifier qu'un entier est nécessairement de la forme $4k, 4k+1, 4k+2,$ ou $4k+3$ et montrer que seules les deux premières formes sont possibles pour son carré.}
		\item\label{q:3}
		Développer $(ax - by)^2 + (ay + bx)^2$ et en déduire que si $m$ et $n$ sont somme de deux carrés, alors le produit $mn$ aussi.
		\item\label{q:4}
		Décomposer 85 en facteurs premiers, écrire chaque facteur comme somme de deux carrés, et en déduire l'écriture de 85 comme somme de deux carrés à l'aide de la question \ref{q:3}.
		\item
		Décomposer $12~240 = a^2 \times b$ avec $a, b$ entiers et $b$ le plus petit possible.
		L'écrire comme somme de deux carrés à l'aide des questions \ref{q:3} et \ref{q:4}.
	\end{enumerate}
}{exe:star2}{
	%À rendre pour possibilité de points bonus à la prochaine évaluation.
	%Ni correction ni points ne seront attribués à un travail suspect ou ne démontrant pas une bonne compréhension de l'exercice.
	
	\underline{Discussion}
	
	%La question 2 deviendra évidente lors de l'étude de l'arithmétique modulaire (programme de Terminale).
	%Essentiellement, les opérations arithmétiques restent cohérentes lorsqu'on étudie uniquement les restes après division euclidienne par 4. 
	%En particulier, le reste du carré est le carré du reste (propriété à montrer !).
	%Les seuls restes possibles pour un carré sont donc $0^2 = 0, 1^2 = 1, 2^2 = 4,$ de reste $0$ et $3^2 = 9$, de reste 1. 
	%On voit donc que les seuls restes 0 et 1 sont possibles.
	
	Nous avons montré que si chaque premier $p$ divisant $n$ et apparaissant dans sa décomposition avec une puissance impaire est somme de deux carrés, alors $n$ est somme de deux carrés (les puissances paires étant déjà des carrés).
	
	Le théorème des deux carrés de Fermat énonce que la direction inverse est également vraie.
	De plus, il énonce que si un premier $p$ n'est pas de la forme $4k+3$ (donc $p=2$ ou $p=4k+1$), alors $p$ est somme de deux carrés.
}
\footnotetext{Pierre de Fermat (1607 - 1665), magistrat et mathématicien français.}


\exemulticols{difficulty=1}{
	Pour quel $N\in\N$ la moyenne de la série ci-contre est-elle 10 ?
}{
		\begin{center}
		\begin{tabular}{|c|c|c|c|c|}\hline
		Valeur   & 0 & 16 & 8 & 12 \\ \hline
		Effectif & 2 & 5 & N & 7  \\ \hline
		\end{tabular}
		\end{center}
}{exe:stats08}{
	Il existe deux façons de procéder : la recherche locale ou la résolution d'une équation.
	
	\underline{Première méthode}
	
	Explorons d'abord la première méthode. Il s'agit de tester une première valeur de $N$, et de décider de la direction (augmenter $N$ ou diminuer $N$) qui permettra de s'approcher de la solution.
	
	Si $N = 0$, alors la moyenne est donnée par
		\[ \frac{0\times2+16\times5+8\times0+12\times7}{2+5+0+7} \approx 11,7. \]
	Comme le but étant d'atteindre 10, il est cohérent d'augmenter le nombre de notes en-dessous de 10 pour baisser la moyenne.
	En augmentant $N$, au augmente donc l'effectif de la note 8, qui va baisser la moyenne.
	
	Si $N=10$, alors la moyenne est donnée par
		\[ \frac{0\times2+16\times5+8\times10+12\times7}{2+5+10+7} \approx 10,2. \]
	Ceci ne suffit pas mais nous évoluons dans la bonne direction !
	
	Si $N=20$, alors la moyenne est donnée par
		\[ \frac{0\times2+16\times5+8\times20+12\times7}{2+5+20+7} \approx 9,5. \]
	Nous sommes allés trop loin : il faut diminuer $N$ pour augmenter la moyenne.
	
	Comme $N \in ]10 ; 20[$, on peut tester $N=15$ et décider d'une direction pour diviser le nombre de valeurs candidates par deux : c'est la dichotomie.

	Si $N=15$, alors la moyenne est donnée par
		\[ \frac{0\times2+16\times5+8\times15+12\times7}{2+5+15+7} \approx 9,79. \]
	Par conséquent, $N \in ]10 ; 15[$.
	
	Si $N=12$, alors la moyenne est donnée par
		\[ \frac{0\times2+16\times5+8\times12+12\times7}{2+5+12+7} = 10, \]
	et nous avons trouvés notre $N$ !

	\underline{Deuxième méthode}

	La contrainte $\overline{X} = 10$ se traduit par l'égalité suivante à résoudre pour $N\in\N$.
	
		\begin{align*}
			\overline{X} &= 10 \\
			\frac{0\times2+16\times5+8\times N+12\times7}{2+5+N+7} &= 10 \\
			0  + 80 + 8\cdot N + 84 &= 10 \cdot(14 + N) \\
			164 + 8 \cdot N &= 140 + 10 \cdot N \\
			164-140 &= 2 \cdot N \\
			N &= 12		
		\end{align*}

	$12$ est bien un entier naturel, donc $N=12$ est l'unique solution.
	Remarquons d'ailleurs que nousa vons démontré l'unicité de la solution en plus de l'avoir obtenue.
}


% TODO: fix leur/leurs ça me saoule ça je dois apprendre
\exe{,difficulty=1}{
	Un élève reçoit deux notes, 7 et 14, et son méchant professeur souhaite modifier leur coefficient afin que l'élève ait une moyenne inférieure à 10.
	\begin{enumerate}
		\item
		Calculer la moyenne des notes lorsque leurs coefficients sont égaux.
		\item
		Que dire sur les coefficients pour que la moyenne pondérée soit inférieure à 10 ?
		\item
		Trouver deux coefficients tels que la moyenne soit \emph{exactement} de 10.
		
		\emph{Indication : supposer que la somme des coefficients est égale à 1.}
		% est-ce que ça va vraiment aider, ça ?
	\end{enumerate}
}{exe:moyenne-ponderee}{
	\begin{enumerate}
		\item
		%D'après l'exercice \ref{exe:3}, lorsque les coefficients sont égaux, ils peuvent être supposés tous les deux égaux à 1.
		Lorsque les coefficients sont égaux, ils peuvent être supposés tous les deux égaux à 1.
		On obtient donc une moyenne classique $\frac{7+14}{2} = 10,5$.
		\item
		Pour obtenir une moyenne pondérée inférieure à 10, il faut augmenter le poids de la note la plus basse.
		Le coefficient associé à 7 doit donc être supérieur à celui associé à 14.
		\item
		Posons $c>0$ le coefficent associé à 7.
		Nous pouvons diviser les coefficients par leur somme sans changer la moyenne pondérée, et donc supposer que la somme des coefficients vaut 1. 
		Le coefficient associé à 14 est donc $1-c$, et la moyenne pondérée vaut
			\[ \frac{7\times c + 14\times(1-c)}{1} = 7c + 14(1-c). \]
		La contrainte de l'énoncé nous permet donc d'obtenir
			\begin{align*}
				7c+14(1-c) &= 10, \\
				7c + 14 - 14c &= 10, \\
				-7c &= -4, \\
				c &= \frac47.
			\end{align*}
		Les coefficients sont donc $\frac47$ et $\frac37$.
		Notons qu'une stratégie de recherche par dichotomie comme à l'exercice \ref{exe:stats08} n'aurait pas abouti car les nombres $\frac47$ et $\frac37$ ne sont pas décimaux : leur développement décimal est infini.
	\end{enumerate}
}


\exe{, difficulty=2}{
	On dit d'une série qu'elle est symétrique autour de 0 si chaque valeur $x$ de la série peut être associée à une unique valeur $-x$ de la série.
	Ainsi $(-1 ; 0 ; 2 ; -2 ; 1)$ est symétrique autour de 0 mais $(1 ; 0 ; -1 ; -1)$ ne l'est pas.
	\begin{enumerate}
		\item
		Montrer qu'une série symétrique autour de 0 a une moyenne nulle.
		\item
		Soit $n\geq1$ un entier naturel.
		À l'aide de la linéarité de la moyenne, montrer que la série statistique $(0 ; 1 ; 2 ; \hdots ; n-1 ; n)$ a pour moyenne $\frac{n}2$.
		\item
		En déduire que
			\[ 1 + 2 + 3 + \dots + (n-1) + n = \dfrac{n(n+1)}2 \]
		\item
		Calculer sans calculatrice la valeur de la somme $1 + 2 + 3 + \hdots+ 9~999 + 10~000$.
	\end{enumerate}
}{exe:somme-triangle}{
%	À rendre pour possibilité de points bonus...

	Indication pour la première partie :
%	montrer que la série est symétrique par rapport à $\frac{n}2$.
%	Chaque valeur $\frac{n}2 - k$ a une sœur $\frac{n}2 + k$.
%	Sommer les valeurs sœurs ensemble pour simplifier le calcul de la moyenne.
	calculer la nouvelle moyenne après avoir soustrait $\frac{n}2$ à chaque valeur de la série statistique et utiliser la linéarité de la moyenne.
}



%%%%%%%%%%%%%%%%%%%%
%%%%% Samoyeau below %%%%%
%%%%%%%%%%%%%%%%%%%%




\newpage
\addcontentsline{toc}{section}{Solutions}
\fancyhead[C]{\textbf{Chapitre \thechapter~— Solutions}}

\shipoutAnswer



